%*************************************
% ภาคผนวก ช.
%*************************************
\vspace{-0.5in}

\appendixChapterTitle{ข}{การเริ่มต้นระบบแบบ Development ในสภาพแวดล้อมหลายรีโพซิทอรี}
\section{การเริ่มต้นบริการประกอบจากรีโพซิทอรี Yuuka}
\hspace*{1.3em}
รีโพซิทอรี Yuuka ทำหน้าที่เก็บไฟล์สำหรับโครงสร้างพื้นฐานของระบบ เช่น Docker Compose, Kubernetes manifests, ตัวอย่าง environment และเครื่องมือ observability ดังนั้นในสภาพแวดล้อมพัฒนา ผู้พัฒนาควรเริ่มต้นบริการประกอบจากรีโพซิทอรีนี้ก่อน เพื่อให้ PostgreSQL, Redis, RustFS และเครื่องมือสำหรับติดตามระบบพร้อมใช้งาน

\begin{figure}[htbp]
\centering
\adjustbox{frame, width=0.95\textwidth}{
    \begin{minipage}{0.9\textwidth}
        \texttt{\footnotesize
        cd yuuka\\
        docker compose -f docker/docker-compose.yaml up -d
        }
    \end{minipage}
}
\caption{\fontSixTeen{การเริ่มต้นบริการประกอบจากรีโพซิทอรี Yuuka}}
\label{figG:StartInfra}
\end{figure}

\section{การเริ่มต้นบริการหลักของระบบ}
\hspace*{1.3em}
หลังจากบริการประกอบพร้อมใช้งานแล้ว ให้เริ่มต้นบริการหลักของระบบใน terminal แยกกัน เพื่อให้สามารถตรวจสอบ log และควบคุมการทำงานของแต่ละบริการได้อย่างชัดเจน

\begin{mycustomenum2}
    \item เริ่ม Momoi จากรีโพซิทอรี \textbf{momoi}
    \item เริ่ม Yuzu จากรีโพซิทอรี \textbf{yuzu}
    \item เริ่ม Midori จากรีโพซิทอรี \textbf{midori}
\end{mycustomenum2}

\begin{figure}[htbp]
\centering
\adjustbox{frame, width=0.9\textwidth}{
    \begin{minipage}{0.85\textwidth}
        \texttt{\footnotesize
        cd momoi\\
        bun run dev
        }
    \end{minipage}
}
\caption{\fontSixTeen{การเริ่มต้น Momoi ในโหมดพัฒนา}}
\label{figG:StartMomoi}
\end{figure}

\begin{figure}[htbp]
\centering
\adjustbox{frame, width=0.9\textwidth}{
    \begin{minipage}{0.85\textwidth}
        \texttt{\footnotesize
        cd yuzu\\
        bun run dev
        }
    \end{minipage}
}
\caption{\fontSixTeen{การเริ่มต้น Yuzu ในโหมดพัฒนา}}
\label{figG:StartYuzu}
\end{figure}

\begin{figure}[htbp]
\centering
\adjustbox{frame, width=0.9\textwidth}{
    \begin{minipage}{0.85\textwidth}
        \texttt{\footnotesize
        cd midori\\
        bun run dev
        }
    \end{minipage}
}
\caption{\fontSixTeen{การเริ่มต้น Midori ในโหมดพัฒนา}}
\label{figG:StartMidori}
\end{figure}

\hspace*{1.3em}
สำหรับ Satsuki โดยทั่วไปจะไม่จำเป็นต้องรันในเครื่องพัฒนาเดียวกันทุกครั้ง เพราะบริการนี้ถูกออกแบบให้ทำงานบนเครื่อง Nginx ภายนอกที่ต้องมีสิทธิ์เขียนไฟล์ configuration และ \textbf{authorized\_keys} จริง หากต้องการทดสอบส่วนนี้ ผู้พัฒนาควรใช้เครื่องทดสอบแยกหรือสภาพแวดล้อมที่เตรียมสิทธิ์ไว้แล้ว

\section{ลำดับการเริ่มต้นที่เหมาะสม}
\hspace*{1.3em}
ลำดับการเริ่มต้นระบบที่เหมาะสมสำหรับการพัฒนา คือ เริ่มจาก Yuuka เพื่อให้บริการประกอบพร้อมก่อน จากนั้นเริ่ม Momoi และ Yuzu แล้วจึงเริ่ม Midori เป็นลำดับสุดท้าย เนื่องจาก Midori ต้องอาศัย API จาก Momoi และ Yuzu ต้องอาศัย Redis รวมถึงฐานข้อมูลในการประมวลผลงานจากคิว

\section{จุดเชื่อมต่อสำคัญในสภาพแวดล้อมพัฒนา}
\hspace*{1.3em}
เมื่อติดตั้งและเริ่มบริการครบแล้ว ผู้พัฒนาสามารถเข้าถึงบริการสำคัญได้จากพอร์ตมาตรฐานดังนี้

\begin{mycustomenum2}
    \item Midori ที่ \url{http://localhost:3000}
    \item Momoi ที่ \url{http://localhost:3000}
    \item Grafana ที่ \url{http://localhost:3002}
    \item RustFS Console ที่ \url{http://localhost:9001}
    \item Prometheus ที่ \url{http://localhost:9090}
    \item Jaeger ที่ \url{http://localhost:16686}
\end{mycustomenum2}

\section{การอัปเดตชนิดข้อมูลระหว่าง frontend และ backend}
\hspace*{1.3em}
เมื่อมีการเปลี่ยนแปลง OpenAPI ของ Momoi ควรเริ่ม Momoi ให้พร้อมก่อน แล้วจึงสั่งสร้างชนิดข้อมูลฝั่ง Midori ใหม่ เพื่อให้ type definitions ของ frontend สอดคล้องกับ API ล่าสุด

\begin{figure}[htbp]
\centering
\adjustbox{frame, width=0.9\textwidth}{
    \begin{minipage}{0.85\textwidth}
        \texttt{\footnotesize
        cd midori\\
        bun run api:gen
        }
    \end{minipage}
}
\caption{\fontSixTeen{การสร้างชนิดข้อมูล API ของ Midori ใหม่จาก OpenAPI}}
\label{figG:ApiGen}
\end{figure}

\section{การแก้ไขปัญหาเบื้องต้นในสภาพแวดล้อมพัฒนา}
\hspace*{1.3em}
แนวทางการตรวจสอบปัญหาที่พบบ่อยในระหว่างการพัฒนามีดังนี้

\begin{mycustomenum2}
    \item \textbf{Port ถูกใช้งานแล้ว}
    \begin{mycustomenum2}
        \item ตรวจสอบ process ที่ใช้งานพอร์ตอยู่
        \item หยุด process ดังกล่าวหรือเปลี่ยนค่าพอร์ตในไฟล์ environment
    \end{mycustomenum2}
    \item \textbf{Docker services ไม่สามารถเริ่มได้}
    \begin{mycustomenum2}
        \item ตรวจสอบสถานะด้วยคำสั่ง \textbf{docker compose -f docker/docker-compose.yaml ps}
        \item ดู log ด้วยคำสั่ง \textbf{docker compose -f docker/docker-compose.yaml logs [service\_name]}
        \item รีสตาร์ทบริการด้วยคำสั่ง \textbf{docker compose -f docker/docker-compose.yaml restart [service\_name]}
    \end{mycustomenum2}
    \item \textbf{Frontend ไม่สามารถเชื่อมต่อ API ได้}
    \begin{mycustomenum2}
        \item ตรวจสอบว่า Momoi ทำงานอยู่ที่พอร์ตที่กำหนดจริง
        \item ตรวจสอบค่า API URL ใน Midori
        \item ตรวจสอบ CORS และ cookie settings ของ Momoi
    \end{mycustomenum2}
    \item \textbf{Yuzu ไม่ประมวลผลงาน}
    \begin{mycustomenum2}
        \item ตรวจสอบ Redis, queue state และ log ของ Yuzu
        \item ตรวจสอบค่า concurrency ของแต่ละคิว
        \item ตรวจสอบ resource lock ที่เกี่ยวข้องกับ \textbf{instanceId} หรือ \textbf{VMID}
    \end{mycustomenum2}
\end{mycustomenum2}
